%----------------------------------
\section{Deliberation on Task 1}
\label{sec:Deliberation}
%----------------------------------

\subsection{Cross-Validation and Sampling Strategy}

This section reconstructs and extends the exact pipeline used in Individual Task 1: Part
1~\cite{nataraja2026task1}, pulled directly from the author's submitted analysis notebook, so that
the cross-validation, learning-curve and bias results below are grounded in the real feature set,
model hyperparameters and split (an 80th-percentile temporal cutoff on Dataset A giving
train~=~23,700 rows and test~=~5,700 rows; a stratified 75/25 split on Dataset B), rather than an
approximation. To check whether these choices delivered an unbiased evaluation, both datasets were
re-evaluated using proper cross-validation schemes --- \texttt{TimeSeriesSplit} for Dataset A and
\texttt{StratifiedKFold} for Dataset B, both from scikit-learn~\cite{pedregosa2011scikit} --- and
compared against the naive alternative (plain random K-fold) that the original report deliberately
avoided. Full code for this comparison is given in Appendix~\ref{sec:Process Documentation}.

\subsubsection*{Dataset A: TimeSeriesSplit vs.\ random K-fold}

A 5-fold random \texttt{KFold} and a 5-fold \texttt{TimeSeriesSplit} were both run on the exact
Gradient Boosting configuration from Task 1 (\texttt{max\_iter=400}, \texttt{learning\_rate=0.06},
\texttt{min\_samples\_leaf=20}) and the exact 12-feature set, including the \texttt{is\_q4}
calendar flag. The random split produced a mean WAPE of 0.668 across folds (individual folds:
0.660, 0.689, 0.674, 0.673, 0.644), while \texttt{TimeSeriesSplit} produced a mean WAPE of 0.634
(folds: 0.579, 0.602, 0.721, 0.666, 0.602).

This is a useful negative result. Random K-fold did not inflate apparent performance here, because
the same-week invoice-count leak already identified and corrected in Task 1: Part 1 was the
dominant leakage source, and it is absent from the current feature set. What random K-fold does
instead is understate performance on some folds and overstate it on others depending on which
weeks land in each validation fold, since the panel is not stationary (weeks vary systematically
with seasonality and promotional timing). \texttt{TimeSeriesSplit}'s middle fold also shows
noticeably higher error (0.721) than its neighbours (0.579, 0.602), which is informative in its
own right: it suggests the relationship between demand and its lagged predictors is not constant
over the two-year window, so a model validated on the wrong slice of time would look better or
worse than it will perform going forward, independently of any leakage. This matters more for an
honest performance estimate than the leakage/no-leakage distinction alone.

A remaining methodological issue surfaced during this re-analysis: the \texttt{price\_rel} feature
(price relative to a product's own median) is computed using each product's median price across
the entire panel, including weeks that fall after the train/test cut-off. This is a mild form of
look-ahead bias --- the median a product settles at over its full lifetime is not knowable at any
earlier forecast time. It did not appear in Task 1: Part 1's discussion of leakage. In a production
pipeline this median should be computed on training-period data only and updated as new weeks
arrive; here it is disclosed rather than corrected retroactively, following Task 1's own principle
of reporting the honest, weaker result.

\subsubsection*{Dataset B: StratifiedKFold vs.\ plain K-fold}

For Dataset B, plain \texttt{KFold} produced validation folds with conversion rates ranging from
14.2\% to 16.7\%, whereas \texttt{StratifiedKFold} held every fold at 15.5\%, matching the
population rate exactly. Mean PR-AUC was similar either way (0.751 vs.\ 0.747), because Gradient
Boosting is reasonably robust to this degree of imbalance drift at this sample size. But the
fold-to-fold PR-AUC spread was slightly tighter under stratification (std 0.015 vs.\ 0.016), and
more importantly, an unlucky plain-KFold split could place a disproportionate share of the
already-rare positive class into one fold, making that fold's estimate unreliable at a 15.5\% base
rate. Task 1's choice of a stratified split for Dataset B (which extends naturally to
\texttt{StratifiedKFold} for cross-validation) is therefore justified by the same logic as WAPE
over MAPE: the metric and the sampling strategy both need to respect the fact that the event of
interest is rare.

\subsection{Learning Curves}

Learning curves were produced for the Gradient Boosting model on both datasets by progressively
increasing the training set size and re-evaluating on the same held-out test set used in Task 1:
Part 1.

\begin{figure}[h]
\centering
\includegraphics[width=0.85\linewidth]{figures/learning_curve_a.png}
\caption{Dataset A --- WAPE on the fixed held-out weeks as training size grows (most recent $k$
weeks used at each step).}
\label{fig:lc_a}
\Description{Line chart showing WAPE falling then plateauing as training set size increases for
Dataset A.}
\end{figure}

WAPE falls from 0.643 at the smallest training size (2,370 rows, roughly two months of the panel)
to 0.589 at the full training set of 23,700 rows, but the descent is noisy rather than smooth
(Figure~\ref{fig:lc_a}): it dips to 0.595 at 18,960 rows, then rises again to 0.609 before falling
to its best value at the full training size. This is not the clean, monotonic improvement a
textbook learning curve shows. More data does help on balance (the full-data WAPE of 0.589 is the
single best point on the curve), but the fluctuations along the way confirm the panel is not
stationary: some additional chunks of historical weeks help the model, and others introduce a
regime that does not transfer as cleanly to the test period, given the non-stationarity already
noted above. The practical implication is that Task 1's WAPE of 0.595--0.607 is a reasonable
estimate of what this feature set can deliver, and further gains are more likely to come from a
wider feature set (promotional flags, competitor price) than from a longer training window alone.

\begin{figure}[h]
\centering
\includegraphics[width=0.85\linewidth]{figures/learning_curve_b.png}
\caption{Dataset B --- PR-AUC on the fixed held-out sessions as (stratified) training size grows.}
\label{fig:lc_b}
\Description{Line chart showing PR-AUC rising steeply then plateauing as training set size
increases for Dataset B.}
\end{figure}

PR-AUC rises steeply from 0.641 at 924 training sessions to about 0.71--0.72 by roughly 4,600
sessions, then plateaus with minor noise (a small dip at 6,472 sessions) through the full training
set of 9,247 (Figure~\ref{fig:lc_b}). The curve has clearly flattened well before the full training
set is used, which means the session classifier is not data-starved: collecting more sessions from
the same population would yield little further gain. Combined with Task 1's finding that
PageValues alone accounts for the majority of permutation importance, this suggests the ceiling on
this task is set by feature richness and the PageValues leakage concern, not by sample size.

\subsection{Bias Analysis}

The assignment's suggested tools (Fairlearn~\cite{bird2020fairlearn}, AI Fairness 360, Google
What-If, LinkedIn Fairness Toolkit) require package installation; Fairlearn could not be installed
in the offline sandboxed analysis environment used for this deliberation (no outbound
package-index access). The bias analysis below instead implements Fairlearn's core method directly
--- disaggregating model error and outcome rates by group (its \texttt{MetricFrame}
approach~\cite{bird2020fairlearn}) --- using pandas and scikit-learn, computing the same group-wise
metrics Fairlearn would report. This choice is justified because neither dataset contains a
protected demographic attribute (race, gender, age, disability, Indigenous status); the available
grouping variables are commercial/behavioural (product revenue tier, visitor type, geographic
region code). A general-purpose disaggregated-metrics approach applies to these groupings just as
well as to demographic ones, and is transparent about exactly which comparison is being made.

\subsubsection*{Dataset A: error by product revenue tier}

\begin{figure}[h]
\centering
\includegraphics[width=0.75\linewidth]{figures/bias_a_tier.png}
\caption{WAPE on held-out weeks, by product revenue tier (thirds of the 300 products, split by
training-period revenue).}
\label{fig:bias_a}
\Description{Bar chart showing WAPE decreasing from the low-revenue tier to the high-revenue
tier.}
\end{figure}

\begin{table}[h]
\centering
\caption{Dataset A: disaggregated forecast error by product revenue tier.}
\label{tab:bias_a}
\small
\begin{tabular}{lrrrr}
\toprule
\textbf{Tier} & \textbf{n (weeks)} & \textbf{MAE} & \textbf{WAPE} & \textbf{Mean demand} \\
\midrule
Low (bottom third)  & 1{,}900 & 133.86 & 0.684 & 195.8 \\
Mid (middle third)  & 1{,}900 & 76.87  & 0.592 & 129.8 \\
High (top third)    & 1{,}900 & 119.93 & 0.510 & 235.3 \\
\bottomrule
\end{tabular}
\end{table}

The model is measurably worse for the lowest-revenue third of products (WAPE 0.684) than for the
highest-revenue third (WAPE 0.510), a 17-point gap (Figure~\ref{fig:bias_a}, Table~\ref{tab:bias_a}).
This is an unbiased-looking model producing a biased business outcome: because training data and
hyperparameter choices are dominated by the volume and revenue patterns of the best-selling lines,
the model implicitly learns their dynamics best. In a pricing and forecasting role, this means
stock and pricing decisions for the long tail of the catalogue --- often the products a specialist
retailer differentiates on --- would be systematically less reliable than for the products that
would sell reasonably well regardless. This is a fairness concern in the commercial rather than
demographic sense, but it is the kind of disparity Carlson and Worrell's broader point about whose
interests a model is optimised for would predict (Part 1): aggregate accuracy metrics (Task 1
reported a single WAPE of 0.595--0.607) conceal exactly this kind of distributional unevenness.

\subsubsection*{Dataset B: error by visitor type and region}

\begin{figure}[h]
\centering
\includegraphics[width=\linewidth]{figures/bias_b_fnr.png}
\caption{False Negative Rate (converting sessions the model misses) by Visitor Type and by Region
code.}
\label{fig:bias_b}
\Description{Two bar charts: false negative rate by visitor type, and false negative rate by
region code.}
\end{figure}

\begin{table}[h]
\centering
\caption{Dataset B: disaggregated error by visitor type.}
\label{tab:bias_b}
\small
\begin{tabular}{lrrrr}
\toprule
\textbf{Visitor type} & \textbf{n} & \textbf{Actual conv.} & \textbf{FNR} & \textbf{FPR} \\
\midrule
New Visitor        & 429   & 25.6\% & 0.245 & 0.050 \\
Returning Visitor  & 2{,}641 & 13.8\% & 0.354 & 0.072 \\
Other              & 13    & 23.1\% & 0.000 & 0.000 \\
\bottomrule
\end{tabular}
\end{table}

Returning visitors who do convert are missed by the model far more often (FNR 35.4\%) than new
visitors who convert (FNR 24.5\%), even though returning visitors make up the large majority of
the test set (Figure~\ref{fig:bias_b}, Table~\ref{tab:bias_b}). A plausible mechanism is that
PageValues and browsing-depth features behave differently for a returning visitor who already
knows what they want and converts quickly, generating a browsing signature closer to a typical
non-converting session. Region-level FNR is more scattered --- from 20.0\% (Regions 5 and 9) to
51.9\% (Region 6) --- but Region 6 also has the smallest positive class among the larger regions
(204 sessions, several with very few converters), so this particular disparity should be treated
as indicative rather than conclusive without more data. The visitor-type gap, backed by a much
larger sample, is the more defensible finding.

Neither disparity was visible in Task 1's single aggregate PR-AUC/F1 figures. If this model
informed which sessions receive a retention offer or a live-chat prompt, converting returning
customers would be disproportionately passed over, which is a materially different business and
fairness outcome from what the headline metrics suggested.

\subsection{Security, Privacy and Ethical Risks at Scale}

\subsubsection*{Security}

Both pipelines were built and validated on static, historical exports. Deployed operationally, the
pricing/forecasting model would need to ingest live transaction and pricing feeds; an adversary
able to inject anomalous transactions (e.g.\ wash sales or scripted bulk orders) could poison the
lagged-demand and rolling-mean features and distort forecasts and any automated reordering or
pricing action downstream. The conversion model, if it scored live sessions, would be exposed to
adversarial probing --- bot traffic engineered to inflate PageValues-adjacent signals could
manipulate which sessions get flagged as high-intent, redirecting promotional spend.

\subsubsection*{Privacy}

Dataset A~\cite{chen2019onlineretailii} includes a Customer ID field not used in the modelling here
but present in the source data; at scale, the same product-week panel could easily be extended to
a customer-week panel, at which point it becomes a behavioural profile of individual purchasing
habits. Dataset B's~\cite{sakar2018shoppers} session-level browsing signals (page depth, exit/bounce
rates, time on page) are similarly re-identifying if ever joined to account-level data or cookies,
even though the released dataset itself is de-identified. Any production deployment would need to
define retention limits, restrict which roles can join this data to identity data, and ensure the
model does not become a de facto tracking layer beyond its stated forecasting purpose.

\subsubsection*{Ethical Risks}

The revenue-tier disparity above means long-tail products and, by extension, the smaller suppliers
or niche lines behind them would receive systematically less reliable forecasting support if this
model informed purchasing or markdown decisions at scale --- an equity concern between product
lines rather than between people, but one with real commercial consequences for whoever depends on
those lines. The visitor-type disparity raises a more direct concern: if conversion scores gated
who receives a discount, a loyalty perk, or a retention call, loyal returning customers who convert
would be under-served relative to new visitors, an outcome opposite to what most retention
strategies intend. Because PageValues partly encodes the outcome it
predicts~\cite{sakar2019realtime} (Task 1, Section~3.3), any high-stakes use of this classifier
beyond the exploratory setting here would need the feature reconstructed to remove that
circularity before deployment, not simply monitored after the fact.

\subsection{Summary}

Re-running Task 1's exact methodology (verified against the author's own submitted notebook, not
an approximation) under proper cross-validation, a learning-curve analysis, and a disaggregated
bias check did not overturn Task 1's headline results, but it added three findings the aggregate
metrics concealed: a subtle look-ahead bias in the \texttt{price\_rel} feature, a learning curve
for Dataset A that improves only noisily rather than monotonically as training data grows, and
materially uneven model quality across product revenue tiers and customer visitor types in both
datasets. None of these change the recommendation already reached in Task 1 --- that neither
dataset alone supports deployment-grade pricing decisions --- but they show that even a
technically sound single-split evaluation can understate how unevenly a model performs across the
population it is meant to serve.
